\section{The two ordinary automorphic seeds over the fixed field}

% BEGIN S69B1
\paragraph{S69B1. The rational ordinary family has unlinked adjoint at every required prime.}
Use the data and the specific primitive holomorphic eigenform $g$ of
P66E1--5, with $k\ge3$, primitive nebentype of conductor $C$ equal to
its level, $\ell\nmid C$, and the prescribed comparison
$\iota:\overline{\mathbb Q}_\ell\simeq\mathbb C$. Write $\pi_u$
for its unitary cuspidal representation. Every finite local parameter
of $\pi_u$ is a sum of two unitary characters with zero monodromy.
For $p\nmid C$, the two roots of the classical Hecke polynomial
have modulus $p^{(k-1)/2}$ by Deligne,
\emph{La conjecture de Weil I}, Thm.~8.2,
doi:10.1007/BF02684373. This statement includes $p=\ell$.
Dividing by $p^{(k-1)/2}$ gives the two unramified unitary parameters.
For $p\mid C$, P66E5 gives a principal series with distinct
inertial characters. Write its parameters $u_1,u_2$ in unitary
normalization. Unitarity identifies their multiset with
$\{\bar u_1^{-1},\bar u_2^{-1}\}$. Each character agrees with its
own conjugate inverse on local units, where it has finite image.
Distinctness there rules out an interchange. Hence each character
is unitary on the entire multiplicative group. These arguments
also establish the assertion for the unitary contragredient.

Let $\rho$ be the actual representation of P66E4--5. At every
$p\ne\ell$, the two eigenvalues of an arithmetic Frobenius lift
in $\rho$ have equal complex absolute value under $\iota$.
At good $p$ this is Deligne's formula. At a bad $p$ use
Langlands--Carayol compatibility for the characteristic-zero
newform, as in P66E5. Its Frobenius-semisimplified Weil--Deligne
parameter is the two characters just obtained, times a common
normalizing character. Monodromy is zero, Frobenius
semisimplification preserves the eigenvalues on Weil elements,
and multiplying a Frobenius lift by inertia changes the two
character values by roots of unity. This proves the assertion
for each choice of place over $p$. Write the two eigenvalues as
$a,b$. The eigenvalues of the trace-zero adjoint are
$a/b,1,b/a$, including if the first matrix has repeated eigenvalues.
Each has complex modulus one. For any two such roots $x,y$ and
any $d\ge1$, $x^d=(p y)^d$ is impossible by absolute values.
Every valuation subring over $p$ admits an arithmetic Frobenius
lift (lift the residue-field Frobenius through the decomposition
group). This lift and exponent one verify the source's
$\mathrm{UnlinkAt}$, in particular whenever the adjoint inertia
image is finite. Along with the very same $\mathrm{EisRep}$ and
$\mathrm{PR}$ from P66E5, this proves the full rational
$\mathrm{PropEGW}$, at the threshold of P66E3 enlarged to at least
three. It supplies both rational families used below.
% END S69B1

% BEGIN S69B2
\paragraph{S69B2. The rational representation for the two-dimensional seed.}
Fix an arbitrary $\mathrm{FieldPairCharacterTripleConditions}$
datum $D$ from $\mathrm{Seed2Alt}$. Identify the fields inside the
specified closure $\Omega$, and write $E=\overline{\mathbb Q}_\ell$,
$G_F=\mathrm{Gal}(\Omega/F)$, $\varepsilon$ for the full cyclotomic
character, and $\Psi=\varepsilon^j\chi'$. Bars denote reduction of
units in $E$. The character $\eta$ cuts out the imaginary
quadratic $K\subset L$, so $\eta(c)=-1$ and $\eta|_{G_L}=1$.
Also $\theta|_{G_L}=1$, and $\theta(c)=1$ since $\theta^\ell=1$
and $\ell$ is odd. By $D$ some $s_0\in G_L$ has
$\bar\Psi(s_0)^2\ne1$.

Use S69B1 with $\phi=(\chi')^2\eta$. Both factors split as
prescribed at $(N,w)$ by $D$, and $\phi(c)=-1$ since $c^2=1$.
In particular $\phi\ne1$. For any $B$, take
$m=j+(\ell-1)(B+k_0+2)$ and $k=2m+1$, where $k_0$ is the
family threshold. Thus $m\ge\max(B,2)$ and the parity and
bound for $k$ hold. Take its $\rho$ and put
\[
  \Xi=\varepsilon^m\chi',\qquad \sigma=\rho\otimes\Xi^{-1}.
\]
The twist is continuous since $\chi'$ factors through
$\mathbb Q(\mu_N)$. The source's determinant and residual
characteristic polynomial for $\rho$ give
\[
 \det\sigma=\eta\theta,\qquad
 \det(X-\bar\sigma(g))=
      (X-\bar\Psi(g)\bar\eta(g))(X-\bar\Psi(g)^{-1}) .
\]
Indeed the residual roots of $\rho(g)$ are $1$ and
$\bar\Xi(g)^2\bar\eta(g)$, and $\bar\Xi=\bar\Psi$ since
$\bar\varepsilon^{\ell-1}=1$. The displayed assertion means
coefficientwise reduction of the characteristic polynomial.

Here is the irreducibility check, including the role of the
particular square condition. In general take such an irreducible
ordinary $\rho$ with $k>1$ and residual roots $(1,\bar\Phi)$.
If its restriction to an open normal subgroup $H$ preserves a
line, choose an inertia element above $\ell$ of infinite
cyclotomic order (the local cyclotomic extensions give open
cyclotomic image). A positive power lies in $H$, with distinct
eigenvalues $\varepsilon^{d(k-1)},1$ by ordinarity, where $d$
is that positive power. Every translate of the line is
$H$-invariant and therefore one of the two eigenlines. Both occur
by global irreducibility. Thus $H$ fixes both and $G_{\mathbb Q}$
acts transitively on their pair. Interchanging elements have
trace zero, so $\bar\Phi=-1$ there. Multiplying by one such
element shows $\bar\Phi=1$ on the subgroup fixing both.
In particular $\bar\Phi|_H=1$ and $\bar\Phi^2=1$ everywhere.
For the present $H=G_L$ (open normal since $L/\mathbb Q$ is
finite abelian), $\bar\Phi(s_0)=\bar\Xi(s_0)^2\ne1$.
Consequently $\rho|_{G_L}$, and also its scalar twist
$\sigma|_{G_L}$, is irreducible.
% END S69B2

% BEGIN S69B3
\paragraph{S69B3. The full rank-two seed, including the all-place witness.}
Put $r=\sigma|_{G_L}$. With
$J=\left(\begin{smallmatrix}0&1\\-1&0\end{smallmatrix}\right)$
and $C_0=\sigma(c)$, set $P=JC_0^{-1}$. The identities
$C_0^2=1$ and $C_0^t J C_0=-J$ give $P^t=P$.
For $s'=csc^{-1}$ and $s\in G_L$, $\det\sigma(s)=1$, so
$r(s)^t P r(s')=P$. The rational residual formula of S69B2
restricts to $(\bar\Psi,\bar\Psi^{-1})$.
At any inertia group away from $\ell$, the two characters of
$\rho$ are simultaneously diagonal in some basis and equal
$(1,(\chi')^2\eta\theta)$ there. On $G_L$ the twist by
$\Xi^{-1}$ gives exactly $((\chi')^{-1},\chi')$ on that
inertia since $\varepsilon$ is unramified. Switch the basis vectors
to get the prescribed order. Above $\ell$ restrict the full
ordinary decomposition-group flag in P66E5 and twist it. Since
$\ell\nmid N$, $\chi'$ is trivial on that inertia. The exponents
change from $(2m,0)$ to $(m,-m)$, and the flag characters reduce
to $(\bar\Psi,\bar\Psi^{-1})$ on the entire decomposition group.
This gives all local conditions for the untwisted $r$.

Now take any $\kappa,\mathrm{ka}$ allowed in $\mathrm{Seed2Alt}$.
The equality in that definition is $\mathrm{ka}=\kappa|_{G_L}$.
By $\mathrm{IsTwist.alg}$, $\kappa$ is the $\ell$-adic realization
of an algebraic Hecke character of $K$. Its norm pullback to $L$
realizes $\mathrm{ka}$ by Artin functoriality. Similarly the
cyclotomic power and finite character of $\Xi^{-1}$ give the
algebraic Hecke character of $\mathbb Q$ realizing that twist.
For the cyclotomic factor the idele norm on a uniformizer has
value $q^{-1}=\varepsilon(\mathrm{Frob}_q^{-1})$.
These comparisons use the source's geometric Frobenius convention.

Start with the regular algebraic cusp attached to $\rho$ in
P66E4. Twist by the indicated character of $\mathbb Q$, base
change via $K$ to $L$, and twist by the indicated character of
$L$. At every embedding the algebraic twists shift both
infinitesimal entries by an integer and base change restricts
the parameter, preserving regular algebraicity. At the real to
complex step the holomorphic discrete-series parameter restricts
to its two distinct complex characters. For cuspidality write
the abelian extension in prime cyclic steps and use cyclic base
change with local compatibility (SX-BC). If a step lost
cuspidality, the preceding cusp would have a nontrivial character
self-twist from that step. Cofinite Frobenius comparison and
Chebotarev would give the same self-twist of the attached
semisimple representation. An intertwiner would be scalar by
irreducibility on the next field (already on $G_L$ by S69B2),
and nonzero scalar intertwiners cannot intertwine a nontrivial
twist. Here the good-place Frobenius comparison propagates at
each step: local base change raises the two parameters to the
residue degree $f$, and $(q^{1/2})^f=(q^f)^{1/2}$. Twisting
multiplies both roots by the Hecke character's geometric
Frobenius value. This proves attachment of the resulting cusp
$\Pi$ to $r'=r\otimes\mathrm{ka}$.

The given twist identity implies
$\mathrm{ka}(s)\mathrm{ka}(csc^{-1})=\varepsilon(s)^{-1}$.
Thus $(r')^c\simeq(r')^\vee\otimes\varepsilon^{-1}$ by $P$.
At good places of norm $q$, dualizing inverts the Galois roots,
and $\varepsilon^{-1}$ on geometric Frobenius contributes $q$.
Consequently $q^{-1/2}$ times the conjugate roots are the inverses
of the original Satake parameters. Strong multiplicity one gives
$\Pi^c\simeq\Pi^\vee$, as required for $\mathrm{AutCsd0}$.

Finally, at every finite place the unitary representation of S69B1
has two unitary support characters. Duality preserves that property
and local base change composes each character with the norm,
even at ramified places. All scalar Hecke twists, including the
normalizing norm twist of P66E4, translate the absolute exponents
by a single real number $t$ independent of the place: the absolute
value of an idele class character factors through the idele norm
by compactness of the norm-one class group. Thus both support
characters have absolute uniformizer value $q^t$, strictly in
the native window $(q^{t-1/2},q^{t+1/2})$, and the component has
toral support. Every regular algebraic cusp attached to $r'$ has
the same good Satake parameters, hence all the same local
components by strong multiplicity one.

To check $\mathrm{PotUnl}$ above $\ell$, twist the flag already
constructed by $\mathrm{ka}$. A power equality between its $i$th
character and $\varepsilon$ times its $j$th character cancels
the common twist. On inertia it would imply
$\varepsilon^{d(h_i-h_j-1)}=1$ for $d\ge1$, where
$(h_1,h_2)=(m,-m)$. The exponent is nonzero. Inertia of any
finite extension of $\mathbb Q_\ell$ has infinite cyclotomic image
because its intersection with the totally ramified cyclotomic tower
is finite. This excludes every power equality, at every valuation
above $\ell$. We obtain $\mathrm{WitA}$ and
$\mathrm{IsAutCsdA}(\iota,r')$. Together with S69B2 this proves
$\mathrm{Sd.Seed2Alt}$ for the specified fields and every $B$.
% END S69B3

% BEGIN S69B4
\paragraph{S69B4. The adjoint seed before automorphic transfer.}
Use the original rank-three data, so $\chi'\ne\eta$, and keep its
$M\ne0$, $2\ell\mid M$, splitting and inertial hypotheses.
This time choose $\phi=\chi'\eta$. The splitting conditions
again multiply. Since $\eta^2=1$, the character $\phi$ is nontrivial.
Take $m=j+(\ell-1)(B+k_0+2)$ for this family's threshold, and set
$k=m+1$. The parity $\chi'(c)=(-1)^j$ and oddness of $\ell$
give $\phi(c)=(-1)^{m+1}$. Take the ordinary $\rho$ of S69B1 at
this $k$, and set $\Phi=\varepsilon^m\chi'\eta$. On $G_L$ we
have $\bar\Phi=\bar\Psi$. In particular $\bar\Phi(s_0)^2\ne1$.
If $\rho$ preserved a line on any open subgroup, it would do so on
that subgroup's open normal core, contradicting the last conclusion
of the invariant-line argument of S69B2. Thus $\rho$ remains
irreducible on every open subgroup.

Let $R=\mathrm{Ad}^0(\rho)|_{G_L}$ on trace-zero two-by-two matrices.
This is the source's $\mathrm{Sym}^2(\rho)\otimes\det(\rho)^{-1}$
in a fixed basis: the alternating determinant pairing identifies
$V^*$ with $V\otimes\det(V)^{-1}$, and the traceless part of
$V\otimes V^*$ corresponds to its symmetric square summand.
The adjoint also stays irreducible on every open subgroup.
Indeed the nondegenerate invariant trace form on
$\mathfrak{sl}_2(E)$ turns an invariant plane into an invariant
orthogonal line. If a line $EX$ were invariant under an open
subgroup $H$ of $G_L$, conjugation by $\rho(H)$ would preserve
the set of eigenlines of the nonzero traceless $X$. A nilpotent
$X$ has one such line. Otherwise $X$ has two distinct eigenlines.
Their stabilizer in $H$ has finite index and is closed by continuity,
hence open. Both cases contradict the open-subgroup irreducibility
of $\rho$, since $G_L$ is open in $G_{\mathbb Q}$. The action of
$R$ is continuous. With $T=\mathrm{Ad}(\rho(c))$, the bilinear
form $P(X,Y)=\mathrm{tr}(X T^{-1}Y)$, where $T^{-1}Y$ means
the image of $Y$ under $T^{-1}$, is nondegenerate and symmetric
because $T$ is an involutive isometry of the trace form. Invariance
under $(R(s),R(csc^{-1}))$ follows from $R(csc^{-1})=TR(s)T^{-1}$.
This gives the required multiplier one.

For $g\in G_L$ let $a,b$ be the roots of the characteristic
polynomial of $\rho(g)$ in $E$. Its integral coefficients and
unit determinant make these roots integral units. Their reductions
are $1,\bar\Phi(g)$ in some order. The adjoint roots
$(a/b,1,b/a)$ therefore give the full coefficientwise residual
polynomial for $(\Psi,1,\Psi^{-1})$. Away from $\ell$ the actual
diagonal inertia of the original $\rho$ has characters
$(1,\chi'\eta\theta)$, which on $G_L$ become $(1,\chi')$.
The off-diagonal matrix units and the traceless diagonal give
simultaneous inertia characters $(\chi',1,(\chi')^{-1})$ in
the appropriate order. Above $\ell$ use the ordinary flag of
$\rho$ with characters $u,v$ in order. In an adapted basis the
upper off-diagonal line, inside the traceless upper triangular
plane, gives an adjoint flag with characters $u/v,1,v/u$.
Their inertia powers are $\varepsilon^m,1,\varepsilon^{-m}$.
On the full decomposition group their reductions, after restriction,
are $\bar\Psi,1,\bar\Psi^{-1}$. All these conclusions hold at
every specified valuation over $L$. As in S69B3 the differences
of the weights $m,0,-m$ exclude $1$ since $m\ge2$. This proves
potential unlinking above $\ell$ after any common scalar twist.
% END S69B4

% BEGIN S69B5
\paragraph{S69B5. The rank-three cusp, its exact normalization and all local windows.}
Let $P_0$ be the base change to $L$ of the \emph{unitary} cusp
$\pi_u^\vee$ of S69B1 for the weight $m+1$ in S69B4. This
base change is cuspidal by the cyclic steps of S69B3: the
algebraically normalized cusp attached to $\rho$ differs from
$\pi_u^\vee$ by one norm twist, and $\rho$ stays irreducible
on every open subgroup. Base change of this twist commutes with
the norm twist over $L$. The same reasoning excludes a
nontrivial Hecke character self-twist of $P_0$. Indeed the central
character forces any such character to have square one, so Artin
reciprocity gives a finite Galois character. At good places
comparison of Frobenius polynomials, after restoring the common
normalization, would give the identical self-twist of $\rho|_{G_L}$
by Chebotarev and semisimplicity. The intertwiner is scalar on
the open kernel, by S69B4, which contradicts a nontrivial twist.

Apply Gelbart--Jacquet, \emph{A relation between automorphic
representations of $\mathrm{GL}(2)$ and $\mathrm{GL}(3)$},
\S\S3 and 9, Thm.~9.3, doi:10.24033/asens.1355. For a unitary
cuspidal representation over a number field with no nontrivial
character self-twist it supplies a cuspidal adjoint lift $P_3$,
with the specified local lift at each place. For principal-series
parameters $u,v$ the local adjoint parameters are $u/v,1,v/u$.
Here every finite place is of this kind: S69B1 gave two unitary
characters, and local base change composes each character with
the local norm. This includes the places above the level $C$,
those above $\ell$, and places ramified in $L/\mathbb Q$. Thus
all three characters of the local support of $P_3$ have modulus
one on uniformizers. They lie in the same open window with
$t=0$ at every finite place; their quotients cannot equal
the absolute value or its inverse. At each complex embedding
the two entries for the restriction of the weight-$m+1$
discrete-series parameter differ by $m$. The adjoint hence has
infinitesimal entries $(m,0,-m)$, all distinct and integral,
the required rank-three regular algebraicity. The unitary cusp
$P_0$ is fixed by CM conjugation since it is a base change from
$\mathbb Q$. The adjoint parameter is self-dual. At good places
these identities give $(P_3)^c\simeq(P_3)^\vee$ by strong
multiplicity one.

To check the normalization, take a good place of $L$ with norm
$q$ and let $a,b$ be the geometric Frobenius eigenvalues of
$\rho|_{G_L}$. Under $\iota$, the two Satake parameters of
$P_0$ differ from $a,b$ by one common scalar
(by P66E4 and the residue-degree formula of S69B3). Thus
the ordinary Satake multiset for $P_3$ is
$\iota\{a/b,1,b/a\}$. Rank-three attachment multiplies it by
$q^{(3-1)/2}=q$. Geometric Frobenius on $\varepsilon^{-1}$ acts
by precisely $q$. Therefore this very $P_3$ attaches
$R\otimes\varepsilon^{-1}$, including its semisimplicity by
S69B4. This verifies $\mathrm{AutCsd0}$. Any other regular
algebraic cusp attached to the same representation has the same
good parameters and hence the same local components by strong
multiplicity one. Together with the all-place window and
S69B4's potential unlinking at every $\ell$-adic place, this gives
the full $\mathrm{WitA}$ and $\mathrm{IsAutCsdA}$ for
$R\otimes\varepsilon^{-1}$.

The $m$, $R$ and symmetric pairing of S69B4 now satisfy
$\mathrm{Sym.UnitaryThreeAutomorphyDense}$ with the original
lower bound $B$ and congruence $m\equiv j\pmod{\ell-1}$.
The direct adjoint lift works for every field and every $M$ in
that statement. The additional splitting and ramification
conditions on $M$ in its alternative unitary-group construction
are consequently allowed throughout. In the linked-end
application, CAF7 supplies the character datum; S69B3 and S69B5
take $B=2$ and give the two actual representations over that
fixed field before any auxiliary space for either rank is chosen.
% END S69B5
